\chapter*{Introduction to Part~II}
\addcontentsline{toc}{chapter}{Introduction to Part~II}
\label{ch:part2-intro}

In this Part we turn from the physics of the inflationary epoch to the present-day dark sector, investigating the nature of dark matter (DM) and dark energy through a combination of cosmological and terrestrial observations.   DM and dark energy together constitute the dominant components of the cosmic energy budget, roughly 95\%, yet their fundamental nature remains among the most pressing open questions in physics. The four chapters that follow attack this problem from complementary directions, unified by a common goal of maximizing the reach of current observations in constraining the properties of the dark sector. This is pursued both by connecting DM production mechanisms to their observational signatures and by developing principled analysis strategies that accurately capture the impact of dark sector physics on cosmic evolution.

The first three chapters focus on DM. In Chapter~\ref{ch:minimal-dm-freezein}, I investigate how the cosmological production of DM via freeze-in is modified when the reheating temperature falls below the DM mass, a regime in which the suppressed production demands larger couplings to match the observed abundance --- bringing the freeze-in scenario within the reach of current direct detection experiments. Chapters~\ref{ch:cmb-leptophilic-dm} and~\ref{ch:pbh-memory-burden} instead explore the power of cosmological data, and in particular the CMB, as a precision probe of exotic energy-injection mechanisms. In Chapter~\ref{ch:cmb-leptophilic-dm}, I show that for light leptophilic DM, CMB anisotropy measurements can constrain decay channels that are partially overlooked in the literature due to their loop suppression, yielding bounds that surpass existing terrestrial limits across a broad range of the parameter space. In Chapter~\ref{ch:pbh-memory-burden}, I turn to a non-particle DM candidate --- primordial black holes --- and test the recently proposed memory-burden mechanism~\cite{Dvali:2018xpy, Dvali:2018ytn, Dvali:2020wft}, which could suppress Hawking evaporation and allow light black holes to survive as DM~\cite{Alexandre:2024nuo, Dvali:2024hsb, Thoss:2024hsr}, against a comprehensive set of cosmological constraints.

Finally, Chapter~\ref{ch:desi-dark-energy} addresses the dark energy sector in the context of the recent claimed evidence from the DESI survey for a time-varying equation of state~\cite{DESI:2024mwx, DESI:2025zgx}. Starting from well-motivated quintessence models, I construct theory-informed priors on the phenomenological parameters that are directly constrained by observations, and use them to critically reassess the robustness of the DESI signal --- illustrating more broadly the importance of principled inference strategies when interpreting marginal detections in cosmological data.
